%=====================================================================
% OSINT Topic 1 — IEEE conference paper
% M.Sc. AI-Enhanced OSINT, SRH University Heidelberg (Prof. Schwarz)
%
% SOURCE OF TRUTH: osint_topic1_fuelcost_pipeline.ipynb
%   Labelling executed at seeds 1234, 5678, 9012 with identical output.
%   + the stationarity patch (ADF, first-differencing) added at Stage 3.
% All reported figures come from a single execution of the notebook:
% the original Stage 3 output (levels) and the stationarity patch
% (ADF + first-differencing) run in the same session.
%=====================================================================
\documentclass[conference]{IEEEtran}
\IEEEoverridecommandlockouts
\usepackage{cite}
\usepackage{amsmath,amssymb,amsfonts}
\usepackage{graphicx}
\usepackage{textcomp}
\usepackage{booktabs}
\usepackage{array}
\usepackage{xcolor}
\def\BibTeX{{\rm B\kern-.05em{\sc i\kern-.025em b}\kern-.08em
    T\kern-.1667em\lower.7ex\hbox{E}\kern-.125emX}}
\begin{document}

\title{Fuel- and Energy-Cost Media Discourse as a Leading Indicator of
German Powertrain Registrations: A Stationarity-Corrected GDELT--KBA Null
Result\\
\thanks{Course project, M.Sc.\ AI-Enhanced OSINT, SRH University
Heidelberg (Prof.\ Dr.\ K.\ D.\ Schwarz), Summer Semester 2026.}}

\author{
\IEEEauthorblockN{Shubhendu Banerjee}
\IEEEauthorblockA{\textit{Dept.\ of Technology and Architecture}\\
\textit{SRH University of Applied Sciences}\\ Heidelberg, Germany\\
ba.shubhendu@gmail.com}
\and
\IEEEauthorblockN{Sacharith Reddy Jambula}
\IEEEauthorblockA{\textit{Dept.\ of Technology and Architecture}\\
\textit{SRH University of Applied Sciences}\\ Heidelberg, Germany\\
sacharithrjambula@gmail.com}
\and
\IEEEauthorblockN{Atharva Shekhar Jadhav}
\IEEEauthorblockA{\textit{Dept.\ of Technology and Architecture}\\
\textit{SRH University of Applied Sciences}\\ Heidelberg, Germany\\
atharvaj567@gmail.com}
\and
\IEEEauthorblockN{Farhan Siddique Mohammed}
\IEEEauthorblockA{\textit{Dept.\ of Technology and Architecture}\\
\textit{SRH University of Applied Sciences}\\ Heidelberg, Germany\\
farhansiddique007890@gmail.com}
}

\maketitle

\begin{abstract}
Open-source discourse is increasingly proposed as an early signal for economic
and consumer behaviour, yet its predictive value for durable-goods markets is
rarely tested with a falsifiable threshold. This paper asks whether the
intensity, thematic composition, and tone of German fuel- and energy-cost news
discourse lead monthly new-vehicle registration shifts across diesel,
battery-electric, and hybrid powertrain categories. We apply a three-stage
open-source intelligence pipeline to 291{,}936 German-language fuel and energy
records from the GDELT Global Knowledge Graph spanning January 2022 to April
2026. Records are embedded with a sentence-transformer model, grouped into 334
discourse communities by weighted community detection, and each community
centroid is labelled once by a grammar-constrained Mistral Nemo 12B classifier
restricted to a seven-value theme taxonomy, so that no record is classified
individually and no label is assigned by hand. Three monthly discourse series
are validated against Kraftfahrt-Bundesamt registration shares over 52 aligned
months. Analysed as levels, the pipeline yields an apparently strong result:
fuel-tag discourse correlates with the hybrid share at $r=-0.698$ and
Granger-predicts it at $p=0.033$. Augmented Dickey-Fuller tests show that the
diesel and hybrid registration shares carry a unit root, so these level results
are spurious. After first-differencing, the smallest p-value across sixteen
tests rises to 0.080, no result approaches the Bonferroni threshold of 0.0031,
and every best-lag bootstrap confidence interval spans the full tested window of
minus six to plus six months. We report a clean, reproducible null, and we
document the level-to-differenced gap as a methodological finding about
discourse-based leading-indicator designs.
\end{abstract}

\begin{IEEEkeywords}
Open-source intelligence, GDELT, community detection, large language models,
Granger causality, stationarity, spurious regression, vehicle registrations,
null result
\end{IEEEkeywords}

%=====================================================================
\section{Introduction}
%=====================================================================
The German passenger-car market is undergoing a structural shift in powertrain
composition. The diesel share of new registrations falls from 0.216 in January
2022 to 0.130 in April 2026, while battery-electric and hybrid vehicles absorb
the released share in an uneven, policy-sensitive pattern. Anticipating these
shifts a few months ahead would carry operational value for manufacturers,
suppliers, grid planners, and policymakers. Because fuel and energy cost is a
plausible economic driver of the choice between combustion and electrified
powertrains, we ask a question that open-source intelligence (OSINT) is well
placed to address. Does public discourse about fuel and energy cost move before
the registration market, or does it merely react to it?

A growing literature reports that media sentiment and search interest carry
predictive information for vehicle and new-energy-vehicle sales
\cite{wijnhoven2017}, \cite{shao2025}, \cite{liu2025}. Most of these studies
operate on curated sentiment corpora and optimise forecast accuracy. Far less
attention is paid to whether a fully open, event-scale source such as the GDELT
Global Knowledge Graph \cite{gdelt2013} carries a signal specific enough to lead
registration behaviour at powertrain granularity, and whether any such lead
survives external validation against an official series the researcher does not
construct.

We address this gap for the German market. Our research question is the
following. To what extent do shifts in the intensity, thematic composition, and
tone of German fuel- and energy-cost discourse in GDELT precede shifts in
Kraftfahrt-Bundesamt (KBA) monthly registration shares for diesel,
battery-electric, and hybrid vehicles, and is the lead time operationally
useful? We fix a falsifiable threshold in advance. The signal is useful only if
the estimated lead time has a bootstrap confidence interval that excludes zero.
If the interval crosses zero, or if no series predicts its registration share on
stationary data, the study reports a null.

This paper makes four contributions.
\begin{itemize}
\item \textbf{A community-then-label OSINT pipeline.} We classify discourse
communities rather than individual records, so a single constrained
language-model call describes each of 334 communities and every record inherits
its community's theme. This removes per-record classifier drift and eliminates
the hand-assigned labels that a stance-mapping design would require.
\item \textbf{A reproducible validation framework.} Every step from stationarity
testing through Granger causality to bootstrap intervals is specified with fixed
parameters and released, providing a template against which a claimed
discourse-based leading indicator can be judged.
\item \textbf{A demonstration that stationarity testing eliminates a false
discovery.} Our own level-based analysis produces a correlation of $-0.698$ and
a Granger p-value of 0.033 for fuel discourse against the hybrid registration
share. We show that both series carry a unit root and that the relationship
disappears under first-differencing, with the same cell rising to $p=0.279$.
\item \textbf{An evidence-based null.} Under the falsifiable
threshold, no discourse series leads any registration share on stationary data.
We report this as a documented negative result, and we note that the knowledge
graph tags no explicit electric-vehicle vocabulary anywhere in this corpus.
\end{itemize}

The remainder of this paper is organised as follows. Section~II reviews related
work. Section~III describes the methodology. Section~IV reports the case study.
Section~V presents results. Section~VI discusses interpretation and threats to
validity, and Section~VII concludes.

%=====================================================================
\section{Related Work}
%=====================================================================
We position this study across four strands: sentiment signals for vehicle
demand, text and news as data for economic forecasting, sequence and OSINT
forecasting systems, and the econometrics of validating a discourse signal.

\subsection{Sentiment Signals for Vehicle Demand}
Consumer sentiment and online attention carry predictive information for car
sales. Combining social-media sentiment with search-trend indices improves
short-horizon forecasts, although raw sentiment often adds little once volume
and search interest are present \cite{wijnhoven2017}. Recent work on new-energy
vehicles shows that media-sentiment features raise the accuracy of
machine-learning sales models \cite{shao2025} and that hybrid
sentiment-and-decomposition architectures lower electric-vehicle forecast error
\cite{liu2025}. These studies establish sentiment as a useful covariate, but
they use curated corpora, target aggregate sales, and optimise accuracy rather
than test a directional lead with a falsifiable threshold. None reports
stationarity handling for the series it correlates.

\subsection{Text and News as Data for Economic Forecasting}
A parallel literature treats text itself as economic data \cite{gentzkow2019}.
Newspaper-derived indicators materially improve forecasts of GDP, inflation, and
unemployment relative to standard baselines \cite{kalamara2022}, and
statistically filtered knowledge graphs built from global news improve forecasts
of industrial production \cite{tilly2021}. Their handling of non-stationarity is
directly relevant to our design. Kalamara et al.\ forecast growth rates rather
than levels and evaluate out of sample, so the trending component of both the
text signal and the outcome is removed before any claim is made
\cite{kalamara2022}, and Tilly and Livan likewise model changes in industrial
production rather than the index itself \cite{tilly2021}. Both therefore avoid
the level-regression problem by construction rather than by explicit unit-root
testing. We make that step explicit instead, testing every series with an
augmented Dickey-Fuller test and reporting the level analysis alongside the
differenced one. This work demonstrates that news text can lead macroeconomic
aggregates, but it does not address a powertrain-resolved consumer-durable
outcome validated against an official registration series.

\subsection{Sequence and OSINT Forecasting Systems}
Transformer architectures now dominate long-horizon time-series forecasting
\cite{informer2021}, and large language models have been applied directly as
zero-shot forecasters \cite{gruver2023}. In the OSINT setting, event-based systems such as
EMBERS forecast civil unrest from open indicators in near real time
\cite{embers2014}, and GDELT tone and theme fields have served as socioeconomic
sensors \cite{gdelt2013}, \cite{bollen2011}. Surveys of the field map its
opportunities and open challenges \cite{osint2020}. These systems are powerful
but data-hungry, and their evaluation rarely states a threshold at which the
signal would be declared absent.

\subsection{Operationalising and Validating a Discourse Signal}
Instruction-tuned language models classify text without task-specific training
data \cite{brown2020}, and grammar-constrained decoding restricts output to a
fixed schema, which makes structured labelling reproducible \cite{willard2023}.
Such classifiers remain reliable only while the target vocabulary is present in
the input, because a constrained model returns a valid label even when the
distinctions its taxonomy encodes are absent, and the failure is silent. This
consideration motivates our decision to label communities rather than records,
and to select a theme taxonomy grounded in the tags the corpus demonstrably
carries. Embeddings for the clustering step follow the Sentence-BERT family
\cite{reimers2019}, and the labelling model is Mistral Nemo \cite{mistralnemo}.

On the validation side the tools are established. Cross-correlation locates
alignment across lags, Granger causality tests whether one series improves
prediction of another beyond its own past \cite{granger1969}, and bootstrap
resampling gives confidence intervals for the estimated lead \cite{efron1993}.
Because both discourse and registration series can trend, spurious-regression
risk is central: Granger and Newbold showed that independent random walks test
significant at conventional levels in the majority of trials, and that
differencing is the standard remedy \cite{grangernewbold1974}. Stationarity is
assessed with the augmented Dickey-Fuller test \cite{dickeyfuller1979}. These
tools have been applied to news-derived series for mobility and macroeconomic
prediction \cite{mobility2021}, and surveys of sentiment-driven econometrics
stress their necessity \cite{algaba2020}.

\subsection{Positioning}
Prior work shows that sentiment predicts vehicle demand, that news text can lead
economic aggregates, and that discourse signals can be validated
econometrically. No prior study combines these into a powertrain-resolved,
stationarity-corrected leading-indicator test for the German market, externally
validated against KBA registrations under a falsifiable
threshold. This paper fills that gap, and reports the size of the error that
omitting the stationarity step would have produced.

%=====================================================================
\section{Methodology}
%=====================================================================
We follow the three-stage OSINT model. Stage one assembles and validates a
fuel-discourse corpus. Stage two embeds records, detects discourse communities,
and labels each community centroid. Stage three builds monthly series and
validates them against registration data with stationarity-aware econometrics.

\subsection{Data Collection and the Corpus Validation Gate}
The corpus is a subset of the GDELT 2.0 Global Knowledge Graph
\cite{gdelt2013}, retrieved through BigQuery and filtered server-side to German
fuel, energy, and automotive coverage. It contains 291{,}936 records covering
January 2022 to April 2026, each carrying a timestamp, a document identifier,
the raw theme list, persons, organizations, and the GDELT tone fields.

Before any pipeline stage runs, the corpus passes a validation gate that
confirms the target signal is present. Of 291{,}936 records, 275{,}880
(94.5\,\%) carry at least one fuel or energy-price theme, with
\texttt{FUELPRICES} on 254{,}040 records, \texttt{ECON\_HEATINGOIL} on
67{,}701, and \texttt{ECON\_GASOLINEPRICE} on 52{,}266. This gate is a direct
response to a diagnostic finding reported in Section~IV, and it is the reason
the discourse construct is fuel and energy cost rather than powertrain stance.

\subsection{Document Representation}
A GKG record carries no article body text, so a text representation must be
constructed from its structured fields. We concatenate human-readable theme
codes with named organizations and persons, which groups records by subject
and actor jointly. Two documented cleaning rules apply. First, GDELT
scaffolding theme families that tokenize into non-topical fragments are removed
by prefix rule, covering \texttt{TAX\_WORLDLANGUAGES}, \texttt{TAX\_ETHNICITY},
\texttt{TAX\_FNCACT}, \texttt{CRISISLEX}, \texttt{UNGP},
\texttt{USPEC\_POLITICS}, and \texttt{MEDICAL} families among others; the
complete prefix list is fixed in the released code.
Second, a document-frequency cutoff removes tokens appearing in more than
40\,\% of records, computed on the final lowercased tokens; this removed 24
tokens. After cleaning, 291{,}291 records carry a usable representation.

\subsection{Embedding and Community Detection}
Records are embedded with all-MiniLM-L6-v2 \cite{reimers2019} as L2-normalised
384-dimensional vectors. Because theme and organization strings repeat heavily,
the 291{,}291 records collapse to 241{,}277 unique strings; identical strings
yield identical embeddings and therefore fall in the same community, so only
unique strings are embedded, each weighted by the number of records it
represents. This makes detection tractable in memory without altering the
partition. Greedy cosine-similarity community detection runs at a threshold of
0.70 with a minimum community size of 250 records, producing 334 communities
covering 203{,}766 records, or 70.0\,\% of the corpus. Each community is
summarised by its centroid and the eight member records nearest to it.

\subsection{Community-Then-Label Classification}
Records are never classified individually. Discovery and categorisation are kept
separate: community detection finds the structure, and the language model maps
each discovered structure to a theme, after which every record inherits its
community's label. This design choice matters for validity. Per-record
classification of 291{,}291 items would invite label drift across calls and
would make the discourse series a function of classifier noise; labelling 334
centroids once makes the mapping auditable in a single table. It also avoids the
alternative failure mode, in which the researcher assigns categories to clusters
by inspection, which would constitute self-constructed ground truth on the
independent variable.

The classifier is Mistral Nemo Instruct 2407 \cite{mistralnemo} in Q5\_K\_M
quantization, run locally with a 4096-token context, thirty-five offloaded
layers, temperature zero, and seed 1234. Output is constrained by a grammar
compiled from a JSON schema \cite{willard2023}, restricting the theme field to
seven values: fuel prices, heating and energy costs, automotive industry, road
transport, energy policy, accidents and safety, and other. The model receives
the eight centroid records and returns a short label and one enum value. No
labelled examples and no fine-tuning are used. Grammar-constrained decoding at
temperature zero makes the labelling deterministic for a fixed prompt, and
repeated execution under seeds 1234, 5678, and 9012 returned identical community
assignments. The reported statistics do not rest on this property in any case:
the fuel-tag share and the fuel-record tone are derived from GDELT theme tags
without reference to the classifier, and an independent implementation that
never loads the model reproduces every statistic in Section~V
(Section~III-H).

\subsection{Discourse Series}
Three monthly series are constructed, ordered from most to least
assumption-free. The \emph{fuel-tag share} is the proportion of monthly records
carrying an explicit fuel or energy-price tag, taken as the union of
\texttt{FUELPRICES}, \texttt{ECON\_GASOLINEPRICE}, and
\texttt{ECON\_HEATINGOIL}; it is independent of both the clustering and the
language model and is therefore the primary series. The \emph{theme shares} are
the per-month proportions of clustered records in each labelled theme. The
\emph{fuel-record tone} is the mean GDELT tone of fuel-tagged records per month,
which becomes more negative as coverage grows more alarmed. Tone is parsed as
missing rather than zero on failure, so that absent tone cannot masquerade as
neutral coverage.

\subsection{External Validation Data}
KBA registration shares are assembled from official products. Diesel counts come
from the FZ~8 and FZ~8.5 releases \cite{kbafz8} and battery-electric and hybrid
counts from the FZ~28 release \cite{kbafz28}. A verified reference month,
January 2022, records 39{,}713 diesel and 20{,}892 battery-electric
registrations out of 184{,}112 total. Shares are the category count divided by
the monthly total, giving three series of 52 months each.

\subsection{Stationarity-Aware Analysis}
Each series is tested with the augmented Dickey-Fuller test using AIC lag
selection \cite{dickeyfuller1979}. Because two registration series carry a unit
root, every series is first-differenced before inference, so that all pairs
receive identical treatment and no result can depend on a shared trend
\cite{grangernewbold1974}. Uniform differencing sacrifices some information from
the already-stationary series, but it provides a consistent stationary
representation across all tested pairs and keeps the results directly
comparable. We report both the level and the differenced analyses, because the
difference between them is itself a finding.

For each pair we compute Pearson cross-correlations across lags from minus six
to plus six months, where a positive lag means discourse leads registrations;
test Granger causality with an SSR-based F test up to four lags
\cite{granger1969}; and estimate a bootstrap confidence interval for the best
lag from 1000 resamples \cite{efron1993}. All econometric tests are computed
with statsmodels \cite{statsmodels}. The 52-month aligned series yield 51
observations after first differencing, which is reflected in the degrees of
freedom of the differenced tests.

We define the inferential testing family explicitly, since the correction
depends on it. Granger tests are run on four discourse-outcome pairs, namely
the fuel-tag share against each of the three registration shares and the
fuel-record tone against the diesel share, at four lags each, giving sixteen
tests per transformation. The theme-share series are not included in this
family: as reported in Section~IV-C they are screened by cross-correlation only
and are treated as descriptive, so they contribute no inferential tests. The
Bonferroni threshold is therefore $0.05/16 = 0.0031$.
Table~\ref{tab:repro} lists the full configuration.

\subsection{Independent Verification}
Because the classifier stage requires a GPU and a quantised 12B model, the
statistical results are additionally reproduced by a second implementation that
depends only on pandas, numpy, and statsmodels. It rebuilds the discourse series
directly from the corpus, re-parses the KBA releases from the published FZ~8,
FZ~8.5, and FZ~28 files, and recomputes the stationarity tests, both Granger
transformations, the cross-correlations, the bootstrap intervals, and the
Johansen test. The two implementations agree on every verdict and to three
decimal places on every statistic, with residual differences arising only from a
record-filtering step that affects 645 of 291{,}936 records. This provides a
check on the reported figures that is independent of the classifier and of the
GPU environment.

\begin{table}[t]
\caption{Reproducibility Configuration}
\label{tab:repro}
\centering \small
\begin{tabular}{@{}ll@{}}
\toprule
\textbf{Element} & \textbf{Value} \\ \midrule
Discourse source & GDELT 2.0 GKG, German fuel/energy \\
Corpus size & 291{,}936 records (2022-01 to 2026-04) \\
Validation gate & 275{,}880 records (94.5\,\%) fuel-tagged \\
Embedding model & all-MiniLM-L6-v2, 384-d, L2-normalised \\
DF cutoff & 0.40 (24 tokens removed) \\
Community threshold & 0.70 cosine \\
Min.\ community size & 250 records \\
Communities & 334, covering 70.0\,\% of records \\
Classifier & Mistral-Nemo-Instruct-2407, Q5\_K\_M \\
Decoding & grammar-constrained JSON, temp 0 \\
Theme enum & 7 values (see Section III-D) \\
Seeds & 1234, 5678, 9012 (identical labels) \\
Validation data & KBA FZ~8, FZ~8.5, FZ~28 (52 months) \\
Stationarity test & augmented Dickey-Fuller, AIC autolag \\
Transformation & uniform first difference \\
Cointegration test & Johansen trace, $k_{ar}$ diff 1 \\
Cross-correlation & Pearson, lags $-6$ to $+6$ \\
Granger test & SSR-based F, maxlag 4 \\
Bootstrap & 1000 resamples \\
Verifier & verify\_values.py (pandas, numpy, \\
 & statsmodels only) \\
Libraries & pandas, sentence-transformers, \\
 & llama-cpp-python, statsmodels, pdfplumber \\
Hardware & single NVIDIA Tesla T4 GPU \\
\bottomrule
\end{tabular}
\end{table}

%=====================================================================
\section{Case Study}
%=====================================================================
\subsection{Why the Construct Is Fuel Cost}
The study was initially designed around powertrain-stance discourse. Diagnostic
probes on the original manufacturer-seeded corpus of 143{,}367 records falsified
that design: the \texttt{ELECTRIC\_VEHICLE} theme matched zero records, and
powertrain vocabulary was effectively absent, with 331 occurrences of
\emph{diesel}, 179 of \emph{bev}, one of \emph{hybrid}, and none of
\emph{battery}. A stance taxonomy applied to that corpus could only have
returned labels unsupported by the underlying text. Fuel and energy-cost tags,
by contrast, are densely and consistently present, and fuel cost is a defensible
economic driver of powertrain mix. The construct was therefore rescoped to what
the instrument measurably contains. Notably, the \texttt{ELECTRIC\_VEHICLE}
theme is absent from the fuel corpus as well, appearing on zero records in all
52 months, so no direct electrification-stance series is available from GDELT at
all.

\subsection{Discourse Corpus Characteristics}
Corpus volume concentrates in the 2022 European energy-price shock and a second
spike in early 2026. Mean tone is negative in every month, ranging from $-1.19$
to $-2.64$, consistent with the crisis framing of energy-cost news. The
fuel-tag share is high and nearly saturated, between 0.837 and 0.987, because
the corpus is already filtered to fuel and energy coverage, so this series
carries limited month-to-month variation.

\subsection{Theme Labelling Outcome}
Table~\ref{tab:themes} reports the distribution of labelled themes across the
203{,}766 clustered records. The distribution is heavily concentrated: the
classifier assigns 87.4\,\% of clustered records to energy policy, with fuel
prices second at 9.9\,\% and the remaining three themes together below 3\,\%.
The concentration indicates that the theme taxonomy provides limited
discrimination within this corpus, and that German fuel and energy discourse in
GDELT is dominated by policy and macroeconomic framing rather than by
consumer-facing price or transport narratives. Theme-share results are
consequently treated as descriptive rather than inferential.

\begin{table}[t]
\caption{Labelled Theme Distribution (Clustered Records)}
\label{tab:themes}
\centering \small
\begin{tabular}{@{}lrr@{}}
\toprule
\textbf{Theme} & \textbf{Records} & \textbf{Share} \\ \midrule
energy policy          & 178{,}063 & 87.4\,\% \\
fuel prices            & 20{,}220  & 9.9\,\% \\
accidents and safety   & 2{,}783   & 1.4\,\% \\
heating and energy costs & 2{,}155 & 1.1\,\% \\
road transport         & 545       & 0.3\,\% \\ \midrule
Total clustered        & 203{,}766 & 100\,\% \\
\bottomrule
\end{tabular}
\end{table}

\subsection{Registration Series}
The KBA diesel share declines steadily from 0.216 in January 2022 to 0.130 in
April 2026, with a trough of 0.108 in December 2022 coinciding with a
battery-electric pull-forward ahead of a subsidy change. The battery-electric
and hybrid shares rise over the window with strong seasonal and policy-driven
spikes. These are the external outcome variables.

%=====================================================================
\section{Results}
%=====================================================================
\subsection{The Level Analysis and What It Appears to Show}
Analysed as levels, the pipeline returns what looks like a positive result.
Fuel-tag discourse correlates with the hybrid registration share at $r=-0.698$
at a two-month lead and with the diesel share at $r=+0.500$ at a three-month
lead. The corresponding Granger tests, reported in the upper rows of
Table~\ref{tab:granger}, place fuel-tag discourse against the hybrid share at
$p=0.033$ at lag two, below the conventional 0.05 threshold. Reported at this
stage, the study would have claimed a two-month leading indicator for hybrid
registrations.

\subsection{Stationarity}
Table~\ref{tab:adf} reports augmented Dickey-Fuller tests on the level series.
The two discourse series and the battery-electric share are stationary, but the
diesel share is emphatically not, at $p=0.998$, and the hybrid share is
non-stationary at $p=0.594$. The diesel share declines monotonically and the
hybrid share rises across the same window, so any level correlation between
discourse and these two outcomes reflects shared trend rather than
co-movement. The level results in Section~V-A are therefore not interpretable
as evidence of a lead.

\begin{table}[t]
\caption{Augmented Dickey-Fuller Tests on Level Series}
\label{tab:adf}
\centering \small
\begin{tabular}{@{}lrrl@{}}
\toprule
\textbf{Series} & \textbf{ADF} & \textbf{p} & \textbf{Verdict} \\ \midrule
Fuel-tag share          & $-3.366$ & 0.012 & stationary \\
Fuel-record tone        & $-2.971$ & 0.038 & stationary \\
KBA diesel share        & $+1.701$ & 0.998 & non-stationary \\
KBA battery-electric share & $-2.962$ & 0.039 & stationary \\
KBA hybrid share        & $-1.376$ & 0.594 & non-stationary \\
\bottomrule
\end{tabular}
\\[2pt]
\footnotesize The five labelled theme-share series were tested under the same
procedure and are stationary at every conventional level ($p \le 0.0006$). They
are omitted here because they carry no inferential tests, as declared in
Section~III-G.
\end{table}

\subsection{Granger Causality Before and After Differencing}
Table~\ref{tab:granger} reports SSR-based F-test p-values for each pair at lags
one to four, in both transformations. The contrast is the central result of this
paper. On levels, the smallest p-value across sixteen tests is 0.033. After
first-differencing, the smallest is 0.080, and the specific cell that carried
the level result, fuel-tag discourse against the hybrid share at lag two, rises
from 0.033 to 0.279. A Bonferroni correction for sixteen tests sets the
threshold at 0.0031, which no cell in either transformation approaches. On
stationary data, no discourse series Granger-predicts any registration share.

\begin{table}[t]
\caption{Granger SSR-Based F-Test p-Values, Levels vs First-Differenced}
\label{tab:granger}
\centering \small
\setlength{\tabcolsep}{4pt}
\begin{tabular}{@{}llrrrr@{}}
\toprule
\textbf{Discourse $\rightarrow$ share} & \textbf{Form} & \textbf{L1} &
\textbf{L2} & \textbf{L3} & \textbf{L4} \\ \midrule
Fuel-tag $\rightarrow$ diesel & levels & 0.132 & 0.250 & 0.110 & 0.181 \\
                              & diff.\ & 0.716 & 0.347 & 0.370 & 0.512 \\
Fuel-tag $\rightarrow$ BEV    & levels & 0.758 & 0.818 & 0.480 & 0.464 \\
                              & diff.\ & 0.514 & 0.183 & 0.244 & 0.372 \\
Fuel-tag $\rightarrow$ hybrid & levels & 0.178 & \textbf{0.033} & 0.067 & 0.067 \\
                              & diff.\ & 0.569 & 0.279 & 0.080 & 0.229 \\
Tone $\rightarrow$ diesel     & levels & 0.259 & 0.713 & 0.798 & 0.532 \\
                              & diff.\ & 0.410 & 0.491 & 0.427 & 0.503 \\
\bottomrule
\end{tabular}
\\[2pt]
\footnotesize Bold marks the only cell below 0.05, which appears in the level
analysis only. Bonferroni threshold for 16 tests is 0.0031. $n=52$ for levels,
$n=51$ after differencing.
\end{table}

\subsection{Cointegration and the Choice of Differencing}
First-differencing is the appropriate treatment only if the non-stationary pairs
do not share a long-run equilibrium relationship, since a vector error-correction
model would otherwise preserve information that differencing discards. We
therefore apply the Johansen trace test to the two pairs whose outcome carries a
unit root. The trace procedure is sequential, so the first hypothesis in the
sequence is decisive. For fuel-tag discourse against the diesel share the trace
statistic at $r \le 0$ is 13.11 against a 95\,\% critical value of 15.49, and
for the hybrid share it is 15.44 against the same critical value. Neither
rejects the null of no cointegrating relationship, so the sequence terminates at
the first step and no cointegrating vector is identified for either pair. The
hybrid statistic falls only marginally below its critical value, and at $n=52$
the finite-sample distribution of the trace statistic warrants caution. On this evidence there is no long-run relationship for a vector
error-correction model to recover, so first-differencing is the appropriate
estimator here rather than a compromise forced by sample size.

\subsection{A Stationary Positive Control}
The battery-electric share provides an internal positive control, because it is
the one registration series already stationary in levels (Table~\ref{tab:adf},
$p=0.039$). A genuine discourse-to-registration lead for this outcome would
therefore be visible in the level analysis without any differencing, and would
not be exposed to the spurious-regression mechanism that invalidates the diesel
and hybrid level results. No such lead appears. Every level p-value for fuel-tag
discourse against the battery-electric share exceeds 0.46
(Table~\ref{tab:granger}), and the strongest level correlation is $+0.144$ at
zero lag (Table~\ref{tab:boot}). The null for this category is therefore not an
artefact of the differencing step, which supports the reading that the absence of
a lead is a property of the discourse signal rather than of the stationarity
correction.

\subsection{Cross-Correlation and Bootstrap Confidence Intervals}
Table~\ref{tab:boot} reports the best lag and its correlation in both
transformations, together with the 95\,\% bootstrap confidence interval on the
differenced series. Differencing reduces every correlation substantially: the
hybrid pairing falls from $-0.698$ to $-0.339$, and the diesel pairing from
$+0.500$ to $+0.255$. For all four pairings the confidence interval spans the
entire tested window, from minus six to plus six months, and none excludes zero.
The falsifiability criterion specified in advance, a lead time whose confidence
interval excludes zero, is not met for any powertrain category.

\begin{table}[t]
\caption{Best-Lag Correlations and Bootstrap Confidence Intervals}
\label{tab:boot}
\centering \small
\setlength{\tabcolsep}{3.5pt}
\begin{tabular}{@{}llrrrrc@{}}
\toprule
& & \multicolumn{2}{c}{\textbf{Levels}} &
\multicolumn{2}{c}{\textbf{Differenced}} & \\
\cmidrule(lr){3-4}\cmidrule(lr){5-6}
\textbf{Discourse} & \textbf{Share} & \textbf{lag} & \textbf{$r$} &
\textbf{lag} & \textbf{$r$} & \textbf{95\,\% CI} \\ \midrule
Fuel-tag & diesel & $+3$ & $+0.500$ & $+3$ & $+0.255$ & $[-6,+6]$ \\
Fuel-tag & BEV    & $0$  & $+0.144$ & $+2$ & $+0.282$ & $[-6,+6]$ \\
Fuel-tag & hybrid & $+2$ & $-0.698$ & $+6$ & $-0.339$ & $[-6,+6]$ \\
Tone     & diesel & $+4$ & $-0.389$ & $+4$ & $-0.231$ & $[-6,+6]$ \\
\bottomrule
\end{tabular}
\end{table}

\subsection{Theme-Share Series}
The labelled theme shares were also screened against the diesel share as a
secondary analysis. On levels the strongest associations are the fuel-prices
theme at a six-month lag ($r=-0.448$), the energy-policy theme at a six-month
lag ($r=+0.400$), and heating and energy costs at a three-month lag
($r=+0.316$). All three peak at negative lags, meaning discourse follows rather
than leads registrations, and given the concentration reported in
Table~\ref{tab:themes} and the stationarity problem established above, we report
these as descriptive only and draw no inference from them.

%=====================================================================
\section{Discussion}
%=====================================================================
\subsection{Interpretation}
The results answer the research question in the negative. Across three discourse
constructions and three powertrain categories, the stationarity-aware Granger
tests show no predictive content at lags one through four, and the
cross-correlation bootstrap intervals identify no stable lead anywhere in the
tested window of minus six to plus six months. German fuel- and energy-cost
discourse in GDELT shows no operationally useful lead over KBA registration
behaviour once the series are made stationary. The Granger
tests do not survive correction, and the bootstrap intervals are as wide as the
search window itself, which is the signature of an estimator that cannot locate
a stable lead. Following Popper, an experiment that could have falsified the
leading-indicator hypothesis and did not produce a signal is evidence
\cite{popper1959}. The threshold fixed in advance makes this a clean,
interpretable null rather than an inconclusive one.

\subsection{The Spurious-Correlation Finding}
The most transferable result of this study is the gap between
Table~\ref{tab:granger}'s two row types. Our pipeline, run in the configuration
a standard lead-lag design would use, produced a correlation of $-0.698$ and a
Granger p-value of 0.033 linking fuel discourse to hybrid registrations. Both
series merely trend across the observation window, one falling as the energy
shock recedes and one rising as electrification proceeds. Under differencing the
same relationship reads $p=0.279$. Had we reported the level result, the
headline finding of this paper would have been a two-month leading indicator
that does not exist.

This has a general implication for discourse-as-signal research. Media attention
to a topic and the adoption curve of a related technology will frequently both
trend over a multi-year window, so the conditions that generate spurious
regression are the normal case rather than an edge case in this literature.
Stationarity testing is therefore not an optional refinement but a precondition
for interpreting any lead-lag claim of this kind. We report both analyses so
that the magnitude of the correction is visible.

\subsection{Why the Null Is Scientifically Valuable}
A negative result under a falsifiable threshold specified in advance is a
contribution rather than a shortfall. Three points support this. First,
research on discourse-based prediction is exposed to publication bias, since
positive findings are reported far more readily than null ones, and a documented
null corrects the record and spares others a redundant search. Second, we
constrain analyst degrees of freedom by fixing the hypothesis and threshold in
advance, reporting the full test grid rather than its smallest cell, and
applying a Bonferroni correction, which together guard against p-hacking. Third,
our own data show that the largest apparent effect in the study is an artefact,
so a discipline that rejects such correlations is exactly what
discourse-based OSINT forecasting requires.

\subsection{Comparison with Prior Work}
Our null appears to contrast with the sentiment-predicts-sales literature
\cite{shao2025}, \cite{liu2025}, but the two are not in direct conflict. That
literature uses curated sentiment corpora and targets aggregate sales, whereas
we test an open event stream at powertrain granularity against official
registration shares, on stationary data, under a falsifiable threshold. Our
finding is consistent with the observation in \cite{wijnhoven2017} that raw
discourse volume and generic sentiment often carry little predictive power once
a strong outcome baseline is present.

\subsection{Threats to Validity and Limitations}
Several limitations qualify the result without overturning it.

\emph{Single market.} The study covers Germany alone through KBA data, so the
null does not generalise to other national markets or regulatory regimes without
replication.

\emph{Observation window.} Fifty-two months is short for Granger analysis, and
51 after differencing. Power is limited, and the window overlaps the post-2022
energy shock, so the single uncorrected differenced result at lag three is best
read as noise.

\emph{Corpus saturation and vocabulary coverage.} The corpus is pre-filtered to
fuel and energy coverage, so the fuel-tag share varies within a narrow band and
volume is driven by the exogenous news cycle rather than consumer intent. GDELT
tags no explicit electric-vehicle vocabulary in this corpus, so the electrified
categories are tested only through fuel-cost proxies.

\emph{Theme concentration.} The classifier assigned 87.4\,\% of clustered
records to a single theme, which limits the discriminative value of the
theme-share series and is why we treat those results as descriptive.

\emph{Classifier variance.} Repeated labelling under three seeds returned
identical assignments, which is expected given deterministic decoding, and the
reported statistics are independent of the classifier in any case. This bounds
only the labelling step. It does not establish that a different embedding model,
similarity threshold, or theme taxonomy would recover the same communities, and
that broader sensitivity analysis is not attempted here.

\emph{Omitted variables.} Bivariate Granger causality does not control for
subsidies, list prices, supply constraints, or seasonality, any of which could
drive both series.

%=====================================================================
\section{Conclusion and Future Work}
%=====================================================================
We asked whether the intensity, thematic composition, and tone of German fuel-
and energy-cost discourse in GDELT lead KBA registration shifts across diesel,
battery-electric, and hybrid powertrain categories, and whether any lead is
operationally useful. The answer is no on all counts. A level-based analysis
appeared to identify a two-month leading indicator for hybrid registrations,
with $r=-0.698$ and $p=0.033$, but augmented Dickey-Fuller tests show the
registration series carry a unit root, and after first-differencing the same
relationship reads $p=0.279$. No differenced result approaches the
Bonferroni-corrected threshold of 0.0031, and every best-lag confidence interval
spans the full search window, so the falsifiable threshold is not met. Beyond
the null, the study contributes a community-then-label pipeline that avoids both
per-record classifier drift and hand-assigned labels, and a documented
demonstration that stationarity testing removes a false discovery that this
design would otherwise have reported.

Several directions follow. Establishing three-run variance bounds across
independent seeds is the immediate next step. A multi-country design using
additional registries, such as the European ACEA series, would test
generalisability. A discourse source carrying explicit powertrain vocabulary,
such as specialist automotive or trade press, social-media streams, or policy
documents, could supply the electrification signal GDELT lacks. On the
modelling side, cointegration testing with a longer window, Bayesian structural
time series, and formal causal-inference designs would strengthen
identification beyond bivariate Granger tests, and a multimodal extension
combining discourse with price and policy series is a natural next step. In each
case the validation framework introduced here provides the baseline against
which a claimed leading indicator must be judged.

%=====================================================================
\begin{thebibliography}{00}

\bibitem{wijnhoven2017} F. Wijnhoven and O. Plant, ``Sentiment analysis and
Google Trends data for predicting car sales,'' in \emph{Proc. 38th Int. Conf.
Information Systems (ICIS)}, Seoul, South Korea, 2017.

\bibitem{shao2025} J. Shao, J. Hong, M. Wang, and X. Wang, ``New energy
vehicles sales forecasting using machine learning: The role of media
sentiment,'' \emph{Computers \& Industrial Engineering}, vol.~201, 2025,
Art.~110928. [Online]. Available: https://doi.org/10.1016/j.cie.2025.110928

\bibitem{liu2025} J. Liu, H. Pan, R. Luo, H. Chen, Z. Tao, and Z. Wu, ``An
electric vehicle sales hybrid forecasting method based on improved sentiment
analysis model and secondary decomposition,'' \emph{Engineering Applications of
Artificial Intelligence}, vol.~150, 2025, Art.~110561. [Online]. Available:
https://doi.org/10.1016/j.engappai.2025.110561

\bibitem{gdelt2013} K. Leetaru and P. A. Schrodt, ``GDELT: Global data on
events, location, and tone, 1979--2012,'' in \emph{Proc. Int. Studies Assoc.
Annu. Conv.}, San Francisco, CA, USA, 2013.

\bibitem{gentzkow2019} M. Gentzkow, B. Kelly, and M. Taddy, ``Text as data,''
\emph{Journal of Economic Literature}, vol.~57, no.~3, pp.~535--574, 2019.
[Online]. Available: https://doi.org/10.1257/jel.20181020

\bibitem{kalamara2022} E. Kalamara, A. Turrell, C. Redl, G. Kapetanios, and
S. Kapadia, ``Making text count: Economic forecasting using newspaper text,''
\emph{Journal of Applied Econometrics}, vol.~37, no.~5, pp.~896--919, 2022.
[Online]. Available: https://doi.org/10.1002/jae.2907

\bibitem{tilly2021} S. Tilly and G. Livan, ``Macroeconomic forecasting with
statistically validated knowledge graphs,'' arXiv:2104.10457, 2021. [Online].
Available: https://arxiv.org/abs/2104.10457

\bibitem{informer2021} H. Zhou, S. Zhang, J. Peng, S. Zhang, J. Li, H. Xiong,
and W. Zhang, ``Informer: Beyond efficient transformer for long sequence
time-series forecasting,'' in \emph{Proc. AAAI Conf. Artificial Intelligence},
vol.~35, no.~12, 2021, pp.~11106--11115. [Online]. Available:
https://doi.org/10.1609/aaai.v35i12.17325

\bibitem{gruver2023} N. Gruver, M. Finzi, S. Qiu, and A. G. Wilson, ``Large
language models are zero-shot time series forecasters,'' in \emph{Adv. Neural
Inf. Process. Syst. (NeurIPS)}, vol.~36, 2023, arXiv:2310.07820. [Online].
Available: https://arxiv.org/abs/2310.07820

\bibitem{embers2014} N. Ramakrishnan \emph{et al.}, ```Beating the news' with
EMBERS: Forecasting civil unrest using open source indicators,'' in \emph{Proc.
20th ACM SIGKDD Int. Conf. Knowledge Discovery and Data Mining (KDD)}, New
York, NY, USA, 2014, pp.~1799--1808.

\bibitem{bollen2011} J. Bollen, H. Mao, and X. Zeng, ``Twitter mood predicts
the stock market,'' \emph{Journal of Computational Science}, vol.~2, no.~1,
pp.~1--8, 2011. [Online]. Available:
https://doi.org/10.1016/j.jocs.2010.12.007

\bibitem{osint2020} J. Pastor-Galindo, P. Nespoli, F. G\'omez M\'armol, and
G. Mart\'inez P\'erez, ``The not yet exploited goldmine of OSINT:
Opportunities, open challenges and future trends,'' \emph{IEEE Access},
vol.~8, pp.~10282--10304, 2020. [Online]. Available:
https://doi.org/10.1109/ACCESS.2020.2965257

\bibitem{brown2020} T. B. Brown \emph{et al.}, ``Language models are few-shot
learners,'' in \emph{Adv. Neural Inf. Process. Syst. (NeurIPS)}, vol.~33, 2020,
pp.~1877--1901.

\bibitem{willard2023} B. T. Willard and R. Louf, ``Efficient guided generation
for large language models,'' arXiv:2307.09702, 2023. [Online]. Available:
https://arxiv.org/abs/2307.09702

\bibitem{reimers2019} N. Reimers and I. Gurevych, ``Sentence-BERT: Sentence
embeddings using Siamese BERT-networks,'' in \emph{Proc. 2019 Conf. Empirical
Methods in Natural Language Processing (EMNLP-IJCNLP)}, Hong Kong, 2019,
pp.~3982--3992. [Online]. Available: https://doi.org/10.18653/v1/D19-1410

\bibitem{mistralnemo} Mistral AI and NVIDIA, ``Mistral NeMo,'' Mistral AI Blog,
Jul.~2024. [Online]. Available: https://mistral.ai/news/mistral-nemo (accessed
Aug.~7, 2026).

\bibitem{granger1969} C. W. J. Granger, ``Investigating causal relations by
econometric models and cross-spectral methods,'' \emph{Econometrica}, vol.~37,
no.~3, pp.~424--438, 1969. [Online]. Available:
https://doi.org/10.2307/1912791

\bibitem{efron1993} B. Efron and R. J. Tibshirani, \emph{An Introduction to the
Bootstrap}. New York, NY, USA: Chapman \& Hall, 1993.

\bibitem{grangernewbold1974} C. W. J. Granger and P. Newbold, ``Spurious
regressions in econometrics,'' \emph{Journal of Econometrics}, vol.~2, no.~2,
pp.~111--120, 1974. [Online]. Available:
https://doi.org/10.1016/0304-4076(74)90034-7

\bibitem{dickeyfuller1979} D. A. Dickey and W. A. Fuller, ``Distribution of the
estimators for autoregressive time series with a unit root,'' \emph{Journal of
the American Statistical Association}, vol.~74, no.~366, pp.~427--431, 1979.
[Online]. Available: https://doi.org/10.2307/2286348

\bibitem{mobility2021} S. Sinha and M. Chakraborty, ``Causal analysis and
prediction of human mobility in the U.S. during the COVID-19 pandemic,''
arXiv:2111.12272, 2021. [Online]. Available: https://arxiv.org/abs/2111.12272

\bibitem{algaba2020} A. Algaba, D. Ardia, K. Bluteau, S. Borms, and K. Boudt,
``Econometrics meets sentiment: An overview of methodology and applications,''
\emph{Journal of Economic Surveys}, vol.~34, no.~3, pp.~512--547, 2020.
[Online]. Available: https://doi.org/10.1111/joes.12370

\bibitem{kbafz8} Kraftfahrt-Bundesamt, ``Neuzulassungen von Kraftfahrzeugen
nach Kraftstoffarten (FZ~8, FZ~8.5),'' Flensburg, Germany. [Online].
Available: https://www.kba.de (accessed Aug.~7, 2026).

\bibitem{kbafz28} Kraftfahrt-Bundesamt, ``Neuzulassungen von
Personenkraftwagen nach ausgew\"ahlten Merkmalen (FZ~28),'' Flensburg,
Germany. [Online]. Available: https://www.kba.de (accessed Aug.~7, 2026).

\bibitem{statsmodels} S. Seabold and J. Perktold, ``Statsmodels: Econometric
and statistical modeling with Python,'' in \emph{Proc. 9th Python in Science
Conf.}, Austin, TX, USA, 2010, pp.~92--96.

\bibitem{popper1959} K. R. Popper, \emph{The Logic of Scientific Discovery}.
London, UK: Hutchinson, 1959.

\end{thebibliography}

\end{document}